\documentclass[11pt]{article}

\usepackage[a4paper,margin=1in]{geometry}
\usepackage[T1]{fontenc}
\usepackage[utf8]{inputenc}
\usepackage{lmodern}
\usepackage{microtype}
\usepackage{amsmath,amssymb,amsthm,mathtools}
\usepackage[hidelinks]{hyperref}
\usepackage{enumitem}

\title{Poincare Inequality as a PL Inequality:\\ Wasserstein, Linearized Wasserstein, and Fixed-Mobility Views}
\author{}
\date{}

\newtheorem{proposition}{Proposition}
\newtheorem{remark}{Remark}
\newtheorem{definition}{Definition}

\DeclareMathOperator{\Ent}{Ent}
\DeclareMathOperator{\Var}{Var}
\newcommand{\dd}{\,\mathrm d}
\newcommand{\R}{\mathbb R}
\newcommand{\W}{W_2}
\newcommand{\F}{\mathcal F}
\newcommand{\E}{\mathcal E}
\newcommand{\D}{\mathcal D}
\newcommand{\I}{\mathcal I}
\newcommand{\dist}{\mathrm{dist}}
\newcommand{\nablapi}{\nabla_{\!\pi}\!\cdot}

\begin{document}
\maketitle

\begin{abstract}
This note explains in which sense the Poincare inequality is, and is not, a Wasserstein Polyak--Lojasiewicz inequality. Log-Sobolev is the genuine $W_2$-PL inequality for relative entropy. Poincare is its linearization at equilibrium, and is exactly a PL inequality for the chi-square energy in the fixed-mobility, or $H^{-1}(\pi)$, transport geometry. If one insists on the usual $W_2$ geometry, Poincare can also be repackaged as a $W_2$-PL inequality for a particular $3/2$-power divergence, but this representation is less canonical.
\end{abstract}

\section{Setup and PL convention}

Let
\[
        \pi(\dd x)=Z^{-1}e^{-V(x)}\dd x
\]
be a smooth probability measure on $\R^d$ or on a smooth Riemannian manifold. Write
\[
        \rho=h\pi, \qquad \int h\dd\pi=1,
\]
and denote by
\[
        L=\Delta-\nabla V\cdot \nabla
\]
the reversible generator, so that
\[
        \int f(-Lg)\dd\pi=\int \nabla f\cdot \nabla g\dd\pi.
\]

For a metric $\mathsf d$ and an energy $\E$, the PL inequality with constant $\kappa>0$ is
\begin{equation}
        2\kappa\bigl(\E-\E_\star\bigr)
        \le |\partial \E|_{\mathsf d}^2 .
        \tag{PL}
\end{equation}
With this convention, a gradient flow satisfying the energy identity
\[
        \frac{\dd}{\dd t}\E(x_t)=-|\partial \E|_{\mathsf d}^2(x_t)
\]
has
\[
        \E(x_t)-\E_\star
        \le e^{-2\kappa t}\bigl(\E(x_0)-\E_\star\bigr).
\]

\section{Log-Sobolev is the genuine Wasserstein--PL inequality}

The relative entropy is
\[
        \Ent_\pi(\rho)
        =\int h\log h\dd\pi.
\]
Its $W_2$ slope is the Fisher information
\begin{equation}
        |\partial \Ent_\pi|_{W_2}^2(\rho)
        =\int \left|\nabla \log h\right|^2h\dd\pi
        =: I_\pi(h).
\end{equation}
Therefore the logarithmic Sobolev inequality
\begin{equation}
        2\alpha\Ent_\pi(h)\le I_\pi(h)
        \tag{LSI}
\end{equation}
is exactly the $W_2$-PL inequality for the entropy, with PL constant $\alpha$.
The associated $W_2$-gradient flow is the Fokker--Planck equation
\[
        \partial_t h_t=Lh_t,
\]
and LSI gives
\[
        \Ent_\pi(h_t)\le e^{-2\alpha t}\Ent_\pi(h_0).
\]

\section{Poincare is the linearized entropy Wasserstein--PL inequality}

The Poincare inequality is
\begin{equation}
        \lambda\int g^2\dd\pi
        \le
        \int |\nabla g|^2\dd\pi,
        \qquad \int g\dd\pi=0.
        \tag{PI}
\end{equation}
It is obtained by linearizing LSI around equilibrium. Indeed, take
\[
        h_\varepsilon=1+\varepsilon g,
        \qquad \int g\dd\pi=0.
\]
Then, as $\varepsilon\to0$,
\[
        \Ent_\pi(h_\varepsilon)
        =\frac{\varepsilon^2}{2}\int g^2\dd\pi+O(\varepsilon^3),
        \qquad
        I_\pi(h_\varepsilon)
        =\varepsilon^2\int |\nabla g|^2\dd\pi+O(\varepsilon^3).
\]
Thus LSI with constant $\alpha$ implies PI with spectral-gap constant at least $\alpha$:
\[
        \alpha\int g^2\dd\pi
        \le
        \int |\nabla g|^2\dd\pi.
\]
Equivalently,
\[
        \alpha_{\rm LSI}\le \lambda_{\rm PI}.
\]
So Poincare is the quadratic, near-equilibrium shadow of the nonlinear entropy/Fisher-information $W_2$-PL inequality.

\section{Why Poincare is not simply $W_2$-PL for chi-square}

Consider the chi-square energy
\[
        \E_2(\rho)
        :=\frac12\chi^2(\rho\mid\pi)
        =\frac12\int (h-1)^2\dd\pi.
\]
For the usual Wasserstein metric, the slope is not the Poincare Dirichlet form. Since
\[
        \frac{\delta \E_2}{\delta \rho}=h-1,
\]
Otto's formula gives
\begin{equation}
        |\partial \E_2|_{W_2}^2(\rho)
        =\int h|\nabla h|^2\dd\pi.
        \tag{$W_2$ slope of $\chi^2$}
\end{equation}
Hence the $W_2$-PL inequality for $\E_2$ would be
\begin{equation}
        \kappa\int (h-1)^2\dd\pi
        \le
        \int h|\nabla h|^2\dd\pi.
        \tag{$W_2$-PL-$\chi^2$}
\end{equation}
This is not the usual Poincare inequality, which has mobility $1$:
\[
        \lambda\int (h-1)^2\dd\pi
        \le
        \int |\nabla h|^2\dd\pi.
\]
The distinction is the mobility. The usual $W_2$ geometry uses the moving mobility $\rho=h\pi$; Poincare uses the fixed equilibrium mobility $\pi$. Near equilibrium, $h\simeq1$, so the two slopes agree to first order.

\section{The exact PL interpretation: fixed-mobility transport / $H^{-1}(\pi)$}

Freeze the Wasserstein mobility at equilibrium. For a mean-zero perturbation $u$, define
\begin{equation}
        \|u\|_{-1,\pi}^2
        :=
        \inf\left\{
        \int |v|^2\dd\pi:
        u=-\nablapi v
        \right\},
        \qquad
        \nablapi v:=\frac1\pi\nabla\cdot(\pi v).
\end{equation}
Equivalently, if $-L\psi=u$, then
\[
        \|u\|_{-1,\pi}^2
        =\int |\nabla\psi|^2\dd\pi
        =\langle u,(-L)^{-1}u\rangle_{L^2(\pi)}.
\]
This is the tangent norm obtained by linearizing $W_2$ at $\pi$.

The corresponding fixed-mobility dynamic distance is
\begin{equation}
        \mathsf d_{-1,\pi}^2(h_0,h_1)
        :=
        \inf_{h_t,v_t}
        \int_0^1\int |v_t|^2\dd\pi\dd t,
        \qquad
        \partial_t h_t+\nablapi v_t=0.
\end{equation}
This is Hilbertian: essentially
\[
        \mathsf d_{-1,\pi}(h_0,h_1)
        =\|h_1-h_0\|_{-1,\pi}.
\]
It is the linearized $W_2$ metric, not the full nonlinear $W_2$ distance.

\begin{proposition}[Poincare as exact $H^{-1}(\pi)$-PL]
For
\[
        \E_2(h)=\frac12\int (h-1)^2\dd\pi,
\]
one has
\begin{equation}
        |\partial \E_2|_{-1,\pi}^2(h)
        =\int |\nabla h|^2\dd\pi.
\end{equation}
Consequently, PI is exactly
\begin{equation}
        2\lambda \E_2(h)
        \le
        |\partial \E_2|_{-1,\pi}^2(h).
        \tag{$H^{-1}$-PL for $\chi^2$}
\end{equation}
\end{proposition}

\begin{proof}
A tangent vector has the form
\[
        \dot h=-\nablapi v.
\]
Then
\[
        \frac{\dd}{\dd s}\E_2(h+s\dot h)\bigg|_{s=0}
        =\int (h-1)\dot h\dd\pi
        =\int \nabla h\cdot v\dd\pi.
\]
By Cauchy--Schwarz in $L^2(\pi)$, the squared dual norm of this differential is
\[
        \int |\nabla h|^2\dd\pi.
\]
Thus the $H^{-1}(\pi)$ slope of $\E_2$ is precisely the Poincare Dirichlet form.
\end{proof}

The associated gradient flow is
\[
        \partial_t h_t=Lh_t,
\]
and the PL inequality gives the usual spectral-gap decay
\[
        \chi^2(\rho_t\mid\pi)
        \le
        e^{-2\lambda t}\chi^2(\rho_0\mid\pi).
\]
Thus the clean interpretation is:
\[
\boxed{\text{Poincare is the PL inequality for }\chi^2/2\text{ in the linearized Wasserstein geometry.}}
\]

\section{A genuine $W_2$-PL inequality generated by Poincare}

There is also a way to encode the Poincare Dirichlet form as a usual $W_2$ slope, but the energy must be changed.

For an $f$-divergence
\[
        \F_\Phi(\rho)=\int \Phi(h)\dd\pi,
\]
Otto's formula gives
\begin{equation}
        |\partial \F_\Phi|_{W_2}^2(\rho)
        =\int h\bigl(\Phi''(h)\bigr)^2|\nabla h|^2\dd\pi.
\end{equation}
To make this slope equal to the Poincare Dirichlet form, choose
\[
        \Phi''(s)=s^{-1/2}.
\]
A convenient normalization is
\begin{equation}
        \Phi(s)=\frac43s^{3/2}-2s+\frac23,
        \qquad
        \Phi(1)=\Phi'(1)=0.
\end{equation}
Then
\begin{equation}
        |\partial \F_\Phi|_{W_2}^2(\rho)
        =\int |\nabla h|^2\dd\pi.
\end{equation}
Moreover, for every $s\ge0$,
\begin{equation}
        \Phi(s)
        =\frac23(\sqrt{s}-1)^2(2\sqrt{s}+1)
        \le
        \frac23(s-1)^2.
\end{equation}
Therefore
\[
        \F_\Phi(\rho)
        \le
        \frac23\chi^2(\rho\mid\pi).
\]
By PI,
\[
        |\partial \F_\Phi|_{W_2}^2(\rho)
        =\int |\nabla h|^2\dd\pi
        \ge
        \lambda\chi^2(\rho\mid\pi)
        \ge
        \frac{3\lambda}{2}\F_\Phi(\rho).
\]
Thus
\begin{equation}
        2\kappa\F_\Phi(\rho)
        \le
        |\partial \F_\Phi|_{W_2}^2(\rho)
        \qquad\text{with }\kappa=\frac{3\lambda}{4}.
        \tag{$W_2$-PL for a $3/2$ divergence}
\end{equation}
This is a genuine $W_2$-PL inequality, but it is not the canonical spectral-gap formulation: the energy is a specially chosen $3/2$-power divergence, not $\chi^2/2$.

\section{Summary dictionary}

\[
\begin{array}{ccl}
\text{Log-Sobolev} 
&\Longleftrightarrow&
W_2\text{-PL for relative entropy }\Ent_\pi, \\
&& |\partial\Ent_\pi|_{W_2}^2=\displaystyle\int |\nabla\log h|^2h\dd\pi; \\
\text{Poincare}
&\Longleftrightarrow&
\text{linearization of entropy }W_2\text{-PL at }h=1;\\
\text{Poincare}
&\Longleftrightarrow&
H^{-1}(\pi)\text{-PL for }\displaystyle\frac12\chi^2(\rho\mid\pi),\\
&& |\partial(\chi^2/2)|_{-1,\pi}^2=\displaystyle\int |\nabla h|^2\dd\pi;\\
W_2\text{-PL for }\chi^2/2
&\Longleftrightarrow&
\displaystyle \kappa\int(h-1)^2\dd\pi\le \int h|\nabla h|^2\dd\pi,\\
&& \text{a different, moving-mobility inequality;}\\
\text{Poincare}
&\Longrightarrow&
W_2\text{-PL for }\displaystyle\int\left(\frac43h^{3/2}-2h+\frac23\right)\dd\pi.
\end{array}
\]

In one sentence: Poincare is best viewed as the PL inequality obtained from the Wasserstein picture by freezing the mobility $\rho$ at its equilibrium value $\pi$.

\begin{thebibliography}{9}

\bibitem{AGS}
L. Ambrosio, N. Gigli, and G. Savare.
\newblock \emph{Gradient Flows in Metric Spaces and in the Space of Probability Measures}.
\newblock Birkhauser, 2008.

\bibitem{BGL}
D. Bakry, I. Gentil, and M. Ledoux.
\newblock \emph{Analysis and Geometry of Markov Diffusion Operators}.
\newblock Springer, 2014.

\bibitem{OttoVillani}
F. Otto and C. Villani.
\newblock Generalization of an inequality by Talagrand and links with the logarithmic Sobolev inequality.
\newblock \emph{J. Funct. Anal.}, 173(2):361--400, 2000.

\bibitem{Villani}
C. Villani.
\newblock \emph{Optimal Transport: Old and New}.
\newblock Springer, 2009.

\end{thebibliography}

\end{document}
